\section{Filtered cobar preliminaries}
\label{may:preliminaries}

We use the filtered Hopf-algebra method of \citet{May1966}, with
the odd-primary May-type filtration in the form of
\citet[Theorem~3.2.5]{Ravenel2004}, and allow an additional digit-zero
weight needed for the injectivity argument in Section~\ref{may:section}.
For earlier applications of Ravenel's filtration to $V(3)$ and powers
of $\beta_1$, see \citet{LeeRavenel1994}.
We construct this modified filtration below, including its initial page
and convergence. Throughout this section, cohomological degree and weight are denoted by
$(s,w)$.  Unless an Adams--Novikov differential is named explicitly, the
spectral sequences in this section are algebraic spectral sequences of
cobar complexes.  Put
\[
 P=\mathbb F_p[t_1,t_2,\ldots],\qquad
 |t_i|_w=u_i,
 \qquad
 \overline\Delta(t_i)=\sum_{0<l<i}t_l\otimes t_{i-l}^{p^l}.
\]
Reversing all coproducts reverses the cobar conventions and changes none of
the conclusions.  We use the displayed convention.

\begin{lemma}[The filtered cobar calculation]
\label{may:page}
For every integer $M\geq0$, the cobar cohomology of $P$ admits a
multiplicative, degreewise finite spectral sequence whose initial
cohomology page is
\begin{equation}
 \Lambda(h_{i,j}:i\geq1,j\geq0)
 \otimes\mathbb F_p[b_{i,j}:i\geq1,j\geq0].
 \label{may:initial-page}
\end{equation}
The bidegrees and auxiliary filtrations are
\[
\begin{array}{c|c|c}
 & (s,w)& F_M\\ \hline
 h_{i,j}&(1,p^ju_i)&2i-1+M[j=0]\\
 b_{i,j}&(2,p^{j+1}u_i)&p(2i-1+M[j=0]).
\end{array}
\]
Here $[j=0]$ is $1$ or $0$ according as $j=0$ or $j\neq0$.
All differentials lower $F_M$.  If an exterior generator is written
$h_{[a,b)}=h_{b-a,a}$, its first differential is
\begin{equation}
 d_1 h_{[a,b)}=\sum_{a<l<b}h_{[a,l)}h_{[l,b)},
 \qquad d_1b_{i,j}=0.
 \label{may:splitting}
\end{equation}
The same statements hold for the Hopf subalgebra
$P_r=\mathbb F_p[t_1,\ldots,t_r]$, with $i\leq r$ in
\eqref{may:initial-page}.
\end{lemma}

\begin{proof}
Write every exponent in base $p$ and assign a digit factor $t_i^{p^j}$
the filtration $2i-1+M[j=0]$.  For a monomial, add these filtrations with
the digit multiplicities, without carrying.  Carrying $p$ copies of one
digit to the next digit lowers filtration: its decrease is
\[
 p(2i-1+M[j=0])-(2i-1+M[j+1=0])\geq p-1>1.
\]
A nonprimitive coproduct summand of $t_i^{p^j}$ replaces it by
$t_l^{p^j}\otimes t_{i-l}^{p^{j+l}}$ and decreases filtration by exactly
one.  These observations show that this is a multiplicative coalgebra
filtration, and that its associated graded Hopf algebra is the tensor
product of primitive truncated polynomial Hopf algebras
\[
 \bigotimes_{i,j}\mathbb F_p[z_{i,j}]/(z_{i,j}^p).
\]
For one such factor the dual algebra is again a truncated polynomial
algebra: factorial rescaling identifies its basis in exponents less than
$p$ with the ordinary polynomial basis.  The free resolution alternating
multiplication by its generator and its $(p-1)$st power computes its
cohomology as one exterior generator in degree one and one polynomial
generator in degree two.  Their weights are respectively the weight of
$z_{i,j}$ and $p$ times that weight.  Tensoring these resolutions gives
\eqref{may:initial-page}, including the stated auxiliary filtrations.
The exterior square is zero since its cohomological degree is odd and
$p$ is odd.

The filtration-one part of the reduced coproduct gives the splitting
formula for $h$.  For the polynomial generator associated to a digit
$z$, a filtration-one contribution would replace one of its $p$ copies
by a split pair and leave $p-1$ unsplit copies of $z$.  In the
associated-graded cohomology of the $z$ factor, a nonzero multiplicity
has the form $pa+\epsilon$, with $a\geq0$ and
$\epsilon\in\{0,1\}$.  The multiplicity $p-1$ has neither form, so that
component has zero cohomology.  Carries cannot contribute to this first
differential by the strict inequality above.  This proves $d_1b=0$.

In a fixed positive internal weight only finitely many digit factors
occur, every monomial exponent is bounded, and a normalized bar has
positive weight.  Its length is consequently bounded as well.  Thus all
the filtrations under discussion are finite in each weight, and their
spectral sequences converge to the indicated cobar cohomology.  The
construction restricts to $P_r$ because its coproduct uses only the
first $r$ coordinates.
\end{proof}

For later calculations we fix the following notation, always
distinguishing the bare exterior letter $b$ from the indexed polynomial
letters:
\begin{equation}
\begin{gathered}
 x=h_{1,0},\quad a=h_{2,0},\quad c=h_{3,0},\qquad
 y=h_{1,1},\quad b=h_{2,1},\quad z=h_{1,2},\\
 A=b_{1,0},\qquad B=b_{2,0},\qquad C=b_{1,1}.
\end{gathered}
\label{may:low-alphabet}
\end{equation}
An exterior interval $h_{i,j}$ covers positions
$j,\ldots,j+i-1$ in its weight, and a polynomial letter $b_{i,j}$
covers positions $j+1,\ldots,j+i$.  Let $m$ be the number of exterior
letters, let $H_r$ count those covering position $r$, and let
$L=(s-m)/2$ be the number of polynomial letters, counted with
multiplicity.  If only positions $0,1,2$ occur, the alphabet is precisely
\eqref{may:low-alphabet}.  Writing $n_A,n_B,n_C$ for the polynomial
multiplicities gives
\begin{equation}
 w=H_0+p(H_1+L-n_C)+p^2(H_2+n_B+n_C),
 \qquad n_A=L-n_B-n_C.
 \label{may:three-position-weight}
\end{equation}
